\chapter{Dermatitis}
\index{Dermatitis}

\section*{Definition and Overview}
Dermatitis is inflammation of the skin, commonly called eczema, that causes redness, itching, dryness, and sometimes blistering, weeping, or scaling. There are several types, including atopic dermatitis (the common form related to allergy and family history), contact dermatitis (caused by irritants or allergens touching the skin), and seborrhoeic dermatitis (affecting oily areas). Dermatitis can affect people of all ages, is often chronic, and can significantly affect quality of life through itching and sleep disturbance. This chapter explains the types of dermatitis, their causes, and how they are treated.

\section*{Epidemiology}
Dermatitis is one of the most common skin conditions. Atopic dermatitis affects around 10--20\% of children and 1--3\% of adults in developed countries, and its prevalence has increased over recent decades. It often begins in infancy and may settle with age, but many affected people have ongoing dry, sensitive skin. Contact dermatitis is common in workplaces and with cosmetics, and seborrhoeic dermatitis affects around 3--5\% of adults. Dermatitis accounts for a large proportion of visits to GPs and dermatologists.

\section*{Aetiology and Pathophysiology}
\begin{itemize}
    \item \textbf{Atopic dermatitis:} a genetic tendency to a defective skin barrier and overactive immune response, triggered by irritants, allergens, temperature, and stress
    \item \textbf{Irritant contact dermatitis:} direct damage to the skin by substances such as soaps, detergents, solvents, and frequent hand washing
    \item \textbf{Allergic contact dermatitis:} an immune reaction to allergens such as nickel, fragrances, and plants (for example poison ivy)
    \item \textbf{Seborrhoeic dermatitis:} linked to Malassezia yeast on oily skin, affecting the scalp, face, and chest
    \item \textbf{Other types:} stasis dermatitis from poor leg circulation, and nummular (coin-shaped) dermatitis
    \item \textbf{Itch-scratch cycle:} scratching worsens inflammation and damages the skin, perpetuating the condition
\end{itemize}

\section*{Clinical Features}
\begin{itemize}
    \item Intense itching, which is the hallmark symptom
    \item Red, dry, scaly, or thickened patches of skin
    \item Small blisters that may weep and crust, especially in acute contact dermatitis
    \item Cracked, sore, and bleeding skin from scratching
    \item Typical patterns: flexures (elbow and knee creases) in atopic dermatitis, hands and exposed areas in contact dermatitis, and scalp and face in seborrhoeic dermatitis
    \item Flare-ups triggered by irritants, allergens, weather, and stress
    \item Sleep disturbance from night-time itching
\end{itemize}

\section*{Differential Diagnosis}
\begin{itemize}
    \item Psoriasis, which causes thick, silvery scaling plaques
    \item Fungal infections such as tinea, which are scaly and sometimes ring-shaped
    \item Scabies, an intensely itchy infestation with burrows and a rash
    \item Lichen planus and other inflammatory skin conditions
    \item Skin lymphoma, which is rare and persistent
    \item Allergic reactions and drug rashes
\end{itemize}

\section*{When to Seek Medical Care}
See a doctor for persistent, severe, or worsening dermatitis, for symptoms that do not respond to moisturisers and self-care, or when the skin becomes painful, oozing, or crusted, which may indicate infection. A dermatologist may be needed for complex or treatment-resistant cases.

\textbf{Contact your local emergency services for:}
\begin{itemize}
    \item Sudden severe allergic reaction with swelling of the face or lips and difficulty breathing
    \item Widespread blistering rash with fever or feeling very unwell (possible serious reaction)
    \item Signs of severe skin infection spreading rapidly
\end{itemize}

\section*{Diagnosis and Investigations}
\begin{itemize}
    \item \textbf{Clinical examination:} the appearance and distribution of the rash usually indicate the type of dermatitis
    \item \textbf{History:} triggers, exposures, family history, and patterns of flare-ups
    \item \textbf{Patch testing:} used in allergic contact dermatitis to identify specific allergens
    \item \textbf{Skin scrapings:} to exclude fungal infection
    \item \textbf{Allergy testing:} skin-prick or blood tests in selected cases of atopic dermatitis
    \item \textbf{Skin biopsy:} occasionally needed for atypical or persistent rashes
\end{itemize}

\section*{Management}
\begin{itemize}
    \item \textbf{Moisturising:} regular emollients repair the skin barrier and are the foundation of treatment
    \item \textbf{Topical corticosteroids:} reduce inflammation during flare-ups; potency is matched to the site and severity
    \item \textbf{Avoidance of triggers:} identify and avoid irritants, allergens, and aggravating factors
    \item \textbf{Antihistamines:} help with itching, particularly at night
    \item \textbf{Topical calcineurin inhibitors:} steroid-free options for sensitive areas such as the face
    \item \textbf{Wet wraps and bleach baths:} for severe atopic dermatitis under specialist guidance
    \item \textbf{Phototherapy and systemic treatment:} for severe cases, including newer targeted medicines
    \item \textbf{Treating infection:} antibiotics or antivirals when the skin is infected
    \item \textbf{Scalp treatment:} medicated shampoos for seborrhoeic dermatitis
\end{itemize}

\section*{Complications}
\begin{itemize}
    \item Secondary bacterial infection (impetigo or cellulitis) from scratching
    \item Eczema herpeticum, a serious viral infection of eczematous skin
    \item Sleep deprivation and its effects on mood and function
    \item Psychological distress, embarrassment, and social impact
    \item Skin thinning or other side effects from prolonged use of strong steroids
    \item Contact allergies developing from treatment products
\end{itemize}

\section*{Prognosis and Outcomes}
Most dermatitis is well controlled with regular moisturising, trigger avoidance, and short courses of treatment during flare-ups. Atopic dermatitis improves with age in many children, although the skin usually remains dry and sensitive. Contact dermatitis resolves when the causative substance is avoided. Seborrhoeic dermatitis is chronic but manageable with regular treatment. With consistent care, most people achieve good control and normal quality of life.

\section*{Prevention and Self-Care}
\begin{itemize}
    \item Moisturise daily, even when the skin is clear
    \item Use gentle, fragrance-free cleansers and lukewarm water
    \item Pat the skin dry rather than rubbing
    \item Avoid known irritants and allergens, and wear gloves for wet work
    \item Keep nails short to reduce damage from scratching
    \item Wear soft, breathable fabrics such as cotton
    \item Manage stress and maintain a regular sleep routine
    \item Follow the doctor's instructions for treatment and review
\end{itemize}

\section*{Sources and Further Reading}
The following reputable sources provide further information for healthcare students and patients seeking reliable, current advice about dermatitis.

\begin{itemize}
    \item Australasian College of Dermatologists. \textit{Eczema and dermatitis information.} https://www.dermcoll.edu.au/
    \item Eczema Association of Australasia. \textit{Support and information for eczema.} https://www.eczema.org.au/
    \item NHS. \textit{Atopic eczema: symptoms, causes, and treatment.} https://www.nhs.uk/conditions/atopic-eczema/
    \item Mayo Clinic. \textit{Dermatitis: types and treatment.} https://www.mayoclinic.org/diseases-conditions/dermatitis-eczema/
    \item DermNet. \textit{Dermatitis and eczema: information.} https://dermnetnz.org/
    \item MedlinePlus. \textit{Dermatitis.} https://medlineplus.gov/
\end{itemize}
